% 18 November 2026. Prepared for Edward Baker with OpenAI Codex LLM
% assistance. Actual inherited serving model and effort unavailable.
\section{Counted coordinates, complete moments and transport errors}
\label{sec:counted-transfer}

Field convergence becomes an observable estimate only after contour
counts, positive field floors and changing rational denominators have
been paid. We now carry out that transfer for the single polynomial of
Section~\ref{sec:common-time}, following \cite{BakerCM48} and its
assessment \cite{BakerCMReview40}. In particular, the polynomial pair
must agree on overlapping charts; unrelated count-two contours would
not identify one pair throughout the interval.

\subsection{Two contours per chart and one polynomial pair}

Retain the independently established canonical tube
\eqref{eq:domain} on the whole fixed $I$. Continuity of $c$ and
compactness provide finitely many closed subintervals $I_\alpha$ whose
relative interiors cover $I$, with reference times $t_\alpha$, such that
\begin{equation}
 \zeta_\alpha=c(t_\alpha),\qquad
 |c(t)-\zeta_\alpha|\le r_0/32\quad(t\in I_\alpha).
 \label{ctr:chart-drift}
\end{equation}
Use time-independent positively oriented circles
\begin{equation}
 \Gamma_\alpha^-:\ |z-\zeta_\alpha|=3r_0/8,
 \qquad
 \Gamma_\alpha:\ |z-\zeta_\alpha|=r,\quad r=r_0/2.
 \label{ctr:two-circles}
\end{equation}
The finite input constants are
\begin{equation}
 \begin{split}
 h_\alpha&=\min_{I_\alpha\times(\Gamma_\alpha^-\cup\Gamma_\alpha)}|H|>0,\\
 K_{j,\alpha}&=\max_{I_\alpha\times\Gamma_\alpha}|\partial_z^jH|,
                       \qquad 1\le j\le3.
 \end{split}
 \label{ctr:contour-inputs}
\end{equation}
Their existence is not a numerical acquisition of their values.
Let $R$ contain all these contours, and suppose that consistent total
field errors satisfy
\begin{equation}
 \sup_{\substack{t\in I\\|z|\le R}}
 |\partial_z^j(P_N-H)(t,z)|\le\epsilon_j,
                       \qquad 0\le j\le3.
 \label{ctr:field-interface}
\end{equation}
Theorem~\ref{cta:field-theorem} supplies $\epsilon_j=C_jE_N$ for exact
data. The arguments below first use arbitrary supplied $\epsilon_j$;
the finite-data terms are added in \eqref{ctr:total-data-error}.

\begin{lemma}[Counted matching on the entire interval]
\label{ctr:matching}
If $\epsilon_0\le h_\alpha/2$ for every chart, both circles in
\eqref{ctr:two-circles} enclose exactly two polynomial occurrences.
These pairs agree on overlaps and define a single selected pair of
$P_N$ on $I$. Its symmetric coordinates $c_N,D_N$ are real and
continuous, including at a double target. No global strip for $P_N$
or labelled square root of $D_N$ is required.
\end{lemma}

\begin{proof}
Every canonical target lies within $9r_0/32$ of $\zeta_\alpha$;
every exterior occurrence is at least $31r_0/32$ away. Both circles
therefore enclose exactly the same two canonical occurrences and are
uniformly zero-free. Rouch\'e's theorem applies on each because
$|P_N-H|\le h_\alpha/2<|H|$. The polynomial pair inside the outer
circle is consequently already inside its companion inner circle.
If two charts overlap, \eqref{ctr:chart-drift} gives
$|\zeta_\alpha-\zeta_\beta|\le r_0/16$; the $\alpha$ inner disk
lies strictly inside the $\beta$ outer disk because
$3r_0/8+r_0/16=7r_0/16<r_0/2$. Its two polynomial occurrences exhaust
the other disk's count. Hence the pairs agree.

The fixed contour power sums define continuous symmetric coordinates,
locally holomorphic in time while the contour stays cleared. Division
by the full counted quadratic gives
\begin{equation}
 H=Q_*A,\qquad P_N=Q_{*,N}A_N,\qquad
 Q_{*,N}=(z-c_N)^2-D_N,
 \label{ctr:factorization}
\end{equation}
with both remaining factors analytic and nonzero in the outer disk.
Since $P_N$ has real coefficients and each disk is conjugation
invariant, its counted pair is conjugation invariant. Thus $c_N,D_N$
are real, whether the pair is real, conjugate nonreal or double.
\end{proof}

\subsection{Quotient, coordinate and denominator errors}

Suppress the chart index temporarily, and write
$B=H_z/H$, $B_N=(P_N)_z/P_N$,
$\langle f\rangle=(2\pi\ii)^{-1}\oint_\Gamma f(z)\dd z$.
The integral satisfies $|\langle f\rangle|\le r\sup_\Gamma|f|$.
Set
\begin{equation}
 b_j=\frac{\epsilon_j}{h-\epsilon_0}
       +\frac{K_j\epsilon_0}{h(h-\epsilon_0)},\qquad
 K_B=K_1/h,\qquad 1\le j\le3.
 \label{ctr:quotient-errors}
\end{equation}

\begin{proposition}[Symmetric-coordinate transfer]
\label{ctr:coordinate-transfer}
Under Lemma~\ref{ctr:matching}, one has $|B_N-B|\le b_1$ and
\begin{equation}
 \begin{split}
 |c_N-c|&\le e_c:=r^2b_1/2,\\
 |D_N-D|&\le e_D:=r^3b_1/2+(r_0/16)e_c+e_c^2.
 \end{split}
 \label{ctr:coordinate-errors}
\end{equation}
On the outer circle put
\begin{equation}
 u_*:=17r_0/32,\qquad q:=161r_0^2/1024,
 \qquad e_Q:=2u_*e_c+e_c^2+e_D.
 \label{ctr:quadratic-errors}
\end{equation}
Then $|Q_*|\ge q$ and $|Q_{*,N}-Q_*|\le e_Q$. If
$e_Q\le q/2$, the two quadratics and every straight interpolation
between them have modulus at least $q/2$ on the outer circle.
\end{proposition}

\begin{proof}
Since $|P_N|\ge h-\epsilon_0$, subtracting the quotients
$\partial_z^jP_N/P_N$ and $\partial_z^jH/H$ proves
\eqref{ctr:quotient-errors}; $j=1$ gives the stated bound for $B$.
Use local coordinates $u=z-\zeta$ and let $I_k=\langle u^kB\rangle$.
The argument principle gives
\[
 c=\zeta+I_1/2,\qquad D=I_2/2-(c-\zeta)^2,
 \qquad |I_{k,N}-I_k|\le r^{k+1}b_1.
\]
Because $|c-\zeta|\le r_0/32$, the difference of the last squares
is at most $(r_0/16)e_c+e_c^2$. This proves
\eqref{ctr:coordinate-errors} without an unnecessary height factor.
On the outer circle,
$15r_0/32\le|z-c|\le17r_0/32$ and $|D|\le r_0^2/16$.
Subtracting the latter allowance from $(15r_0/32)^2$ gives $q$.
Expanding $(z-c_N)^2-(z-c)^2$ gives the bound $e_Q$.
Finally $|Q_*+\theta(Q_{*,N}-Q_*)|\ge q-e_Q$ for $0\le\theta\le1$.
The inner circle similarly has canonical quadratic floor
$57r_0^2/1024$; rational integration below uses only the outer circle.
\end{proof}

The same interface also pays the first time derivatives. Define
\begin{equation}
 A_j=K_j/h,\qquad A_{j,N}=(K_j+\epsilon_j)/(h-\epsilon_0),\qquad
 b_t=b_3+A_{2,N}b_1+A_1b_2.
 \label{ctr:time-quotient}
\end{equation}
The two heat equations and \eqref{eq:log-time} imply
$|(B_N)_t-B_t|\le b_t$. Differentiating the fixed contour formulas gives
\begin{equation}
 \begin{split}
 |c_N'-c'|&\le e_{c'}:=r^2b_t/2,\qquad
 |c'|\le C_{c'}:=r^2(A_3+A_1A_2)/2,\\
 |D_N'-D'|&\le r^3b_t/2+(r_0/16)e_{c'}
                           +2e_c(C_{c'}+e_{c'}).
 \end{split}
 \label{ctr:time-coordinate-errors}
\end{equation}
For the last line, expand the derivative of $(c-\zeta)^2$ and use
$|c_N'|\le C_{c'}+e_{c'}$. These estimates remain valid at either
discriminant zero; they make no Lipschitz assertion for individual roots.

\begin{corollary}[An explicit cutoff threshold]
\label{ctr:threshold}
For exact-data errors $\epsilon_j=C_jE_N$, define on each chart
\begin{equation}
 \begin{gathered}
 b=2C_1/h+2K_1C_0/h^2,\qquad A_c=r^2b/2,\\
 A_D=r^3b/2+(r_0/16)A_c+A_c^2,\qquad
 A_Q=2u_*A_c+A_c^2+A_D,\\
 E_*^{(\alpha)}=\min\{1,h/(2C_0),q/(2A_Q)\}.
 \end{gathered}
 \label{ctr:linear-threshold}
\end{equation}
A zero denominator constant imposes no restriction. If
$E_N\le\min_\alpha E_*^{(\alpha)}$, all preceding count and
denominator tests hold, and $b_1\le bE_N$, $e_c\le A_cE_N$,
$e_D\le A_DE_N$, $e_Q\le A_QE_N$. In particular the explicit
sufficient test is
\begin{equation}
 \frac{2K\{\log(2+N)+1\}}N\le\min_\alpha E_*^{(\alpha)}.
 \label{ctr:cutoff-test}
\end{equation}
\end{corollary}

\begin{proof}
The second restriction in \eqref{ctr:linear-threshold} gives
$h-\epsilon_0\ge h/2$, hence the asserted linear bound for $b_1$.
Use $E_N^2\le E_N$ in \eqref{ctr:coordinate-errors} and
\eqref{ctr:quadratic-errors}; the third restriction then gives
$e_Q\le q/2$. Equation~\eqref{cta:tail-bound} proves
\eqref{ctr:cutoff-test}.
\end{proof}

\subsection{Complete rational traces and fixed-moment rates}

Define $S_{n,N}$ and $T_{k,N}$ using \emph{every} polynomial occurrence
outside the selected pair, with its actual multiplicity. The canonical
quantities still use every canonical exterior occurrence. Neither
collection is an unmentioned cutoff of an evolved exterior divisor.

\begin{theorem}[Complete exterior-moment transfer]
\label{ctr:moment-transfer}
Assume the count and denominator tests above. Put
\begin{equation}
 U_n=(2/q)^n,\qquad V_n=n(2/q)^{n+1}.
 \label{ctr:kernel-constants}
\end{equation}
For every fixed $n\ge1$ and $k\ge2$, the uniform errors on a chart are
\begin{align}
 |S_{n,N}-S_n|&\le e_{S_n}:=r(U_nb_1+K_BV_ne_Q),
 \label{ctr:even-error}\\
 |T_{k,N}-T_k|&\le e_{T_k}:=
 r\{(u_*+e_c)U_kb_1+K_B(e_cU_k+u_*V_ke_Q)\}.
 \label{ctr:odd-error}
\end{align}
For the exact-data threshold of Corollary~\ref{ctr:threshold}, explicit
linear constants are
\begin{equation}
 \begin{split}
 C_{S_n}^{(\alpha)}&=r(U_nb+K_BV_nA_Q),\\
 C_{T_k}^{(\alpha)}&=
 r\{(u_*+A_c)U_kb+K_B(A_cU_k+u_*V_kA_Q)\}.
 \end{split}
 \label{ctr:moment-linear-constants}
\end{equation}
Taking maxima over the finite cover gives errors $C_{S_n}E_N$ and
$C_{T_k}E_N$ on all of $I$, hence $O(\log N/N)$ at each fixed order.
\end{theorem}

\begin{proof}
First identify the traces with the \emph{complete} exterior moments.
For either $R(z)=Q_*^{-n}$ or $R(z)=(z-c)Q_*^{-k}$, all kernel poles
are enclosed targets and $R=O(z^{-2})$ or $O(z^{-3})$ at infinity.
For an exterior occurrence $\rho$, the rational function
$R(z)/(z-\rho)$ has zero residue at infinity and exactly one pole
outside the contour, at $\rho$. Thus
\begin{equation}
 \left\langle\frac{R(z)}{z-\rho}\right\rangle=-R(\rho).
 \label{ctr:individual-residue}
\end{equation}
Also $\langle R\rangle=0$. The combined contribution of the counted
targets is $\langle RQ_*'/Q_*\rangle=0$, because this rational
function has only enclosed poles and decays faster than $1/z$.
This includes a double target; individual simple target residues need
not vanish separately.

For $P_N$, complete partial fractions and
\eqref{ctr:individual-residue} prove the traces immediately.
For $H$, use its complete even order-at-most-one product
\eqref{eq:product}. The regularized logarithmic derivative converges
normally on the fixed circle; its additive constants integrate to zero
against $R$. The inverse-square tail pays this passage, and the
resulting exterior sums converge absolutely. Therefore
\begin{equation}
 S_n=-\langle Q_*^{-n}B\rangle,\qquad
 T_k=-\langle(z-c)Q_*^{-k}B\rangle,
 \label{ctr:complete-traces}
\end{equation}
and the same formulas hold with every object indexed by $N$.
Only individual rational functions were evaluated at infinity; no
uncontrolled outer contour for the full logarithmic derivative has
been discarded. Local nonvanishing of the deflated factor alone would
not identify its derivatives with these complete sums.

The straight interpolation of Proposition~\ref{ctr:coordinate-transfer}
gives
$|Q_{*,N}^{-n}-Q_*^{-n}|\le V_ne_Q$, while each inverse power is
bounded by $U_n$. Subtract the two even traces, using the perturbed
kernel with $B_N-B$ and the kernel difference with $B$, to obtain
\eqref{ctr:even-error}. For the odd kernel, additionally use
$|z-c_N|\le u_*+e_c$ and
\[
 |(z-c_N)Q_{*,N}^{-k}-(z-c)Q_*^{-k}|
       \le e_cU_k+u_*V_ke_Q.
\]
This proves \eqref{ctr:odd-error}. Finally apply
Corollary~\ref{ctr:threshold} and $E_N\le1$ to obtain
\eqref{ctr:moment-linear-constants}.
\end{proof}

Exterior--exterior mergers remain outside the fixed circles and cause
no singularity in these formulas. Internal target mergers are handled
by the quadratic kernel. A target--exterior merger would destroy the
supplied collar and is outside the theorem. The same fixed-contour
formulas, \eqref{ctr:time-coordinate-errors} and differentiation of
their rational kernels also give first-time-jet passage at fixed
moment orders. Constants depend on those orders and the contour floors;
no simultaneous height/order limit follows.

\subsection{Forcing and every-subinterval transport}

Let all errors below denote their maxima over the finite chart cover.
Distinguish the physical forcing from its scaled version:
\begin{equation}
 \begin{gathered}
 M=S_1,\qquad \mathcal F=12S_2+16DS_3,\qquad
 M_t=2M^2-\mathcal F,\qquad D_t=2-4DM,\\
 \sigma=(t-t_0)/\ell^2,\quad \delta=D/\ell^2,\quad
 s_n=\ell^{2n}S_n,\quad v_k=\ell^{2k-1}T_k,\quad m=\ell^2M,\\
 F=\ell^4\mathcal F=12s_2+16\delta s_3.
 \end{gathered}
 \label{ctr:physical-normalized}
\end{equation}
Use identical definitions for the polynomial. If
$A_3^{\rm mom}\ge\sup_I|S_3|$ is a supplied canonical moment or
contour bound, then
\begin{equation}
 e_M=e_{S_1},\qquad
 |\mathcal F_N-\mathcal F|\le e_{\mathcal F}:=
 12e_{S_2}+16\{d\ell^2e_{S_3}
                    +e_D(A_3^{\rm mom}+e_{S_3})\}.
 \label{ctr:forcing-error}
\end{equation}
Indeed write $D_NS_{3,N}-DS_3=D(S_{3,N}-S_3)+(D_N-D)S_{3,N}$.
This retains the product of errors. Consequently
$|m_N-m|\le\eta:=\ell^2e_M$, $|F_N-F|\le\ell^4e_{\mathcal F}$,
$|\delta_N-\delta|\le e_D/\ell^2$, and the even and odd scaled
moment errors are respectively $\ell^{2n}e_{S_n}$ and
$\ell^{2k-1}e_{T_k}$. A proved polynomial forcing inequality may be
transferred using these errors; the error theorem itself supplies no
improved forcing sign.

\begin{theorem}[Every-subinterval multiplier and physical cost]
\label{ctr:transport-theorem}
Put $\Sigma=\Delta/\ell^2$. For every $0\le x\le y\le\Sigma$,
let
\begin{equation}
 W(x,y)=\exp\left(-4\int_x^y m(v)\dd v\right),\qquad
 W_N(x,y)=\exp\left(-4\int_x^y m_N(v)\dd v\right).
 \label{ctr:multipliers}
\end{equation}
Then both multipliers are positive and
\begin{equation}
 \begin{gathered}
 |\log W(x,y)-\log W_N(x,y)|\le4\eta(y-x),\\
 \e^{-4\eta(y-x)}W_N(x,y)\le W(x,y)
                 \le\e^{4\eta(y-x)}W_N(x,y),\\
 |W(x,y)-W_N(x,y)|\le(\e^{4\eta(y-x)}-1)W_N(x,y).
 \end{gathered}
 \label{ctr:multiplier-error}
\end{equation}
The physical discriminant cost is
\begin{equation}
 \mathcal C_D(x,y)=2\ell^2\int_x^yW(v,y)\dd v,
 \qquad
 \mathcal C_{D,N}(x,y)=2\ell^2\int_x^yW_N(v,y)\dd v.
 \label{ctr:physical-cost}
\end{equation}
It satisfies
\begin{equation}
 \begin{split}
 |\mathcal C_D-\mathcal C_{D,N}|
 &\le2\ell^2\int_x^y
             (\e^{4\eta(y-v)}-1)W_N(v,y)\dd v\\
 &\le(\e^{4\eta(y-x)}-1)\mathcal C_{D,N}(x,y).
 \end{split}
 \label{ctr:cost-error}
\end{equation}
At exact fixed inputs these are $O(E_N)=O(\log N/N)$ errors,
uniformly over all subintervals.
\end{theorem}

\begin{proof}
Reality of the two selected pairs and their complete conjugation-invariant
exterior divisors makes $m,m_N$ real. Integrating the uniform bound
$|m-m_N|\le\eta$ proves the logarithmic estimate. Exponentiation gives
the two-sided multiplier inequalities and their absolute consequence.
The variation-of-constants formula for the physical equation is
\begin{equation}
 D(t_0+\ell^2y)
   =W(x,y)D(t_0+\ell^2x)+\mathcal C_D(x,y).
 \label{ctr:variation-constants}
\end{equation}
Here $\partial_\sigma D=2\ell^2-4mD$, explaining the indispensable
factor $\ell^2$ in \eqref{ctr:physical-cost}. Apply the multiplier
bound separately on every $[v,y]$ and integrate to obtain
\eqref{ctr:cost-error}. For fixed inputs, $e_M=O(E_N)$ and $m$ is
bounded; hence $m_N$ and the displayed costs are uniformly bounded
for all sufficiently large $N$. This proves the stated linear rate.
\end{proof}

For an explicit absolute bound, any supplied $|m|\le A_m$ gives
$|m_N|\le A_m+\eta=:A_N^{\rm abs}$. Define
$\Psi(b,h)=\int_0^h\e^{bv}\dd v$, with its continuous value $h$
at $b=0$. With $h=y-x$, \eqref{ctr:cost-error} implies
\begin{equation}
 |\mathcal C_D-\mathcal C_{D,N}|
 \le2\ell^2\{
 \Psi(4(A_N^{\rm abs}+\eta),h)-\Psi(4A_N^{\rm abs},h)\}.
 \label{ctr:absolute-cost-error}
\end{equation}
The normalized $\delta$ cost is $\mathcal C_D/\ell^2$.
In physical time the logarithmic multiplier error is
$4e_M(v-u)$ on every $[u,v]\subset I$.

Independent anchors remain necessary for a sign consumer. If an incoming
anchor has error $\alpha_x$, subtraction in
\eqref{ctr:variation-constants} gives the additional allowance
\[
 W(x,y)\alpha_x+|W-W_N|\,|D_N(t_0+\ell^2x)|
                   +|\mathcal C_D-\mathcal C_{D,N}|.
\]
A supplied terminal lower anchor $A_y$ implies a nonnegative incoming
$D$ if it exceeds $\mathcal C_{D,N}$ plus the upper cost error, because
$W>0$. Approximation acquires neither that anchor nor the tube; equality
and discriminant-zero cases remain included.

\subsection{Finite input errors and the scope of the realization}

\begin{proposition}[Weighted finite-data propagation]
\label{ctr:data-proposition}
Suppose the supplied real-even initial polynomial has coefficient errors
$|\widehat a_{2k}-a_{2k}|\le\xi_k$, $0\le k\le J_N$, and evolves
by the same terminating heat exponential. Then, for $0\le j\le3$,
\begin{equation}
 \epsilon_j^{\rm data}\le
 \sum_{k=0}^{J_N}\xi_k
 \sum_{r=0}^{\lfloor(2k-j)/2\rfloor}
  \frac{\Delta^r}{r!}\frac{(2k)!}{(2k-2r-j)!}R^{2k-2r-j}.
 \label{ctr:coefficient-data-error}
\end{equation}
A negative upper summation limit means zero, and $0^0=1$.
All preceding estimates apply to the supplied polynomial with
\begin{equation}
 \epsilon_j=C_jE_N+\epsilon_j^{\rm data},
 \label{ctr:total-data-error}
\end{equation}
provided the count and denominator tests hold for these total errors.
\end{proposition}

\begin{proof}
Apply $\partial_z^j\exp(-\tau\partial_z^2)$ to each error monomial
$\delta a_{2k}z^{2k}$. Its finite expansion is
\[
 \delta a_{2k}\sum_{r=0}^{\lfloor(2k-j)/2\rfloor}
 \frac{(-\tau)^r}{r!}\frac{(2k)!}{(2k-2r-j)!}z^{2k-2r-j}.
\]
Taking absolute values and using $\tau\le\Delta$, $|z|\le R$
proves \eqref{ctr:coefficient-data-error}. Add the canonical truncation
error from Theorem~\ref{cta:field-theorem}. The resulting polynomial
has the same real symmetry required in Lemma~\ref{ctr:matching}, so
every subsequent estimate applies with the total error interface.
\end{proof}

Location enclosures can be converted into these coefficient budgets,
but only after the finite occurrences have been correctly matched.
If $|\widehat\rho_\nu-\rho_\nu|\le e_\nu<|\rho_\nu|$, put
\begin{equation}
 b_\nu=|\rho_\nu|^{-2},\qquad
 d_\nu=\frac{e_\nu(2|\rho_\nu|+e_\nu)}
               {|\rho_\nu|^2(|\rho_\nu|-e_\nu)^2}.
 \label{ctr:root-data-error}
\end{equation}
The difference-of-inverse-squares identity gives
$|\widehat\rho_\nu^{-2}-\rho_\nu^{-2}|\le d_\nu$.
If the normalization error is at most $e_a$, a sufficient coefficient
allowance is
\begin{equation}
 \xi_k=e_a e_k(b+d)+h_0\{e_k(b+d)-e_k(b)\},
 \label{ctr:symmetric-data-error}
\end{equation}
where $e_k$ is the elementary symmetric polynomial and $e_0=1$.
To see this, expand the two finite products in inverse squares. The
normalization difference contributes the first term; expansion of the
remaining difference and telescoping bounds every mixed term by the
corresponding positive term of $e_k(b+d)-e_k(b)$.

An alternative matched-factor budget, directly suited to Gaussian
propagation, uses $\alpha_\nu=\rho_\nu^{-2}$ and
$w_\nu=\max(|\alpha_\nu|,|\widehat\alpha_\nu|)$:
\begin{equation}
 |\widehat P_N(t_0,z)-P_N(t_0,z)|
 \le\left[e_a+h_0|z|^2\sum_{\nu\le J_N}
                  |\widehat\alpha_\nu-\alpha_\nu|\right]
                    \prod_{\nu\le J_N}(1+|z|^2w_\nu).
 \label{ctr:factor-data-error}
\end{equation}
Finite telescoping proves this inequality. Its whole Gaussian expectation
is an explicit finite polynomial in $R$ and $\sqrt{2\Delta}$, with
Gaussian absolute moments $\mu_0=1$, $\mu_1=\sqrt{2/\pi}$ and
$\mu_j=(j-1)\mu_{j-2}$. Cauchy then pays its spatial jets on a padded
disk. Likewise a supplied initial uncertainty
$\delta f_0(z)=\int_0^\infty\delta\psi_0(u)\cos(zu)\dd u$
has a valid whole-time error budget
\begin{equation}
 V(R,\Delta)=\int_0^\infty|\delta\psi_0(u)|
                          \e^{\Delta u^2+Ru}\dd u<\infty.
 \label{ctr:weighted-source-error}
\end{equation}
The Fubini argument of Theorem~\ref{cta:field-theorem} yields the
evolved source representation with weight $\e^{\tau u^2}$; its
absolute value is bounded by $V$. On a disk padded by $p>0$, the
corresponding order-$j$ jet budget is
$j!p^{-j}V(R+p,\Delta)$. If the initial uncertainty is the sum of a
polynomial error and such a source error, add their budgets. These are
uncertainty budgets for the unchanged target, not a change of physical source.

The degree-dependent weights are part of the result. Fixed unweighted
coefficient accuracy does not imply a uniform conditioning theorem as
$N$ grows. Wrong multiplicities, omitted factors or uncertain cutoff
completeness are not small matched location errors. Exact real even
symmetry must also be retained: a small complex perturbation alone does
not give real $m_N$ or the positive costs used above.

Finally, the direct polynomial method with the same input divisor can
evaluate exactly the same complete signed rational sums as
\eqref{ctr:complete-traces}; it receives every error bound above.
Avoiding exterior root labels is not a proved speed advantage. Generic
finite polynomial controls initialized with a nonzero value at the
origin retain their mergers: after their entire initial divisor is
included, $E_N=0$ and all exact-data errors vanish.
In particular the assessed separated quartic control has simultaneous
target and exterior double zeros while its counted transfer is exact
\cite{BakerCMReview40}. Thus finite realization and approximation do
not provide a strict gap sign. The canonical result is a fixed-input
common-time approximation, distinct from short-head closure, acquired
canonical data, a uniform-height limit or a new Newman bound.
